\section{Methods}
\label{sec:methods}

\subsection{Two audit scopes}
\label{sec:audit-metrics}

We run the directional audit of \cref{eq:dchi_def} at two levels.

\paragraph{Cosmological candidate transfer (main audit).}
Both SPT likelihoods and the completion are placed in a single Cobaya model with one shared Boltzmann calculation, so the two lenses see identical theory spectra at identical parameters. The model rejects any difference between the two sides in theory settings, sampled parameterization, hard prior bounds or completion components. Each lens is fitted by maximizing the sum of its own native log-likelihood and the completion's log-likelihoods within the hard prior bounds, without the separate Cobaya prior density. The native \texttt{candl} priors, including the optical-depth priors, remain part of each lens. The transferred vector $\phi$ contains the seven cosmological parameters, including $\mnu$. In the full CMB+DESI baseline the test lens also re-fits the Planck calibration $A_{\mathrm{planck}}$ and the tSZ and kSZ amplitudes. In the SPT+DESI baseline no parameters are re-fitted. Other SPT nuisance parameters stay at their packaged default values. The joint $\dchi$ includes the completion terms, and we report the SPT and completion contributions separately.

The fits use bounded derivative-free optimization (BOBYQA) from several starting points. The reference fit for the test lens is started from, among other points, the transferred candidate, and the better candidate is retained. All endpoints are re-evaluated with a freshly built model, which reproduces the native likelihoods, spectra and accounting to better than $10^{-10}$ in $\chi^2$. The fits are numerical candidates and are not certified global optima.

\paragraph{Fixed-spectrum nuisance transfer (secondary).}
To isolate nuisance packaging, we also run the audit at fixed theory spectra: the packaged test spectra shipped with each release, not cosmologies fitted to each survey. The tSZ and kSZ amplitudes are shared between the lenses. In addition, a capped set of twelve calibration and foreground parameters of the test lens is re-fitted. The transferred profile starts from the transferred point. The reference fit releases both the shared and the capped parameters and is searched from two starting points: the transferred profile and the lens's own fit. The capped sets are listed in Appendix~\ref{app:reproducibility}. This scope probes how far nuisance re-tuning can absorb the disagreement between the packaged spectra. It does not test cosmological cross-prediction.

\subsection{Exact localization with correlated data}
\label{sec:ledger}

For a Gaussian likelihood with residual $r=d-t(\theta)$ and covariance $C$, the quadratic form is $Q=r^\top C^{-1}r$. A tempting localization inverts each block of $C$ separately and sums the block chi-squares. When blocks are correlated this sum differs from $Q$. For example, with
\[
  C=\begin{pmatrix}1&0.9\\0.9&1\end{pmatrix},\qquad r=\begin{pmatrix}1\\2\end{pmatrix},
\]
the full form is $Q\approx7.37$, while the two one-dimensional block chi-squares sum to $1+4=5$. Instead we allocate the full form to coordinates,
\begin{equation}
  q_i \;=\; r_i\,(C^{-1}r)_i, \qquad \sum_i q_i = Q,
  \label{eq:signed_alloc}
\end{equation}
which shares each cross term $r_iC^{-1}_{ij}r_j$ equally between $i$ and $j$. Summing $q_i$ over any partition into groups $g$ (here spectrum $\times$ multipole band) gives group contributions that add up exactly to $Q$. The audit uses the difference of these allocations between the transferred and reference endpoints,
\begin{equation}
  \Delta Q_g \;=\; \sum_{i\in g}\bigl[q_i(\hat\theta_A)-q_i(\hat\theta_B)\bigr],
  \label{eq:block_dchi}
\end{equation}
where each endpoint uses its own effective covariance. This matters because the SPT 2018 beam covariance depends on the parameters, and a Hartlap correction applies where present. Individual $q_i$ and $\Delta Q_g$ can be negative. They are an accounting of one quadratic form and not independent chi-square variables. Only the full form $Q$ is guaranteed to be nonnegative.

The SPT part of the audit then reconciles exactly with the native likelihood:
\begin{equation}
  \dchi_{\mathrm{SPT}} \;=\; \sum_g\Delta Q_g \;+\; \Delta Q_{\mathrm{unlocalized}} \;+\; \Delta(\text{native prior penalty}) \;+\; \Delta(\log\det C).
  \label{eq:ledger}
\end{equation}
Bins outside the plotted bands are kept as an explicit ``unlocalized'' term. For example, the 2018 TE and EE bins with $300\le\ell<400$ fall below the first band edge. The exact-algebra parts of this ledger (solve, allocation and grouping) are verified in Lean (Appendix~\ref{app:formal}, \texttt{CovarianceLedger}). The numerical reconstruction is checked against the native likelihood in every run.

\subsection{Controls}
\label{sec:common-staging}
We repeat the cosmological audits under three controls, each of which re-runs all four fits:
\begin{itemize}[leftmargin=*, itemsep=2pt]
  \item \textbf{Support cut.} D1 is restricted to TT $750\le\ell\le3000$ and TE/EE $\ell\le3000$, removing its extra coverage, while the 2018 lens is left unchanged. D1 cannot supply the 2018 TE/EE modes with $300\le\ell<400$, and binning and window functions are not matched.
  \item \textbf{Optical-depth prior} (full CMB+DESI). The native SPT Gaussian prior on $\tau$ is removed from both lenses; all other priors and the Planck low-$\ell$ EE likelihood are kept.
  \item \textbf{Boltzmann precision} (full CMB+DESI). Seven CLASS integration, sampling and quadrature settings are raised towards the upstream reference values, with the physical model unchanged.
\end{itemize}

\subsection{Posterior sampling and diagnostics}
\label{sec:mnu-inference}

We sample the seven parameters $\omega_b$, $\omega_c$, $H_0$, $\ln(10^{10}A_s)$, $n_s$, $\tau$ and $\mnu$ with Cobaya's Metropolis sampler. The full CMB+DESI runs also sample $A_{\mathrm{planck}}$ and the tSZ and kSZ amplitudes. The SPT+DESI runs keep all SPT nuisance parameters at their packaged defaults. For CLASS we sample an internal mass parameter, identical to $\mnu$, and pass $m_{\mathrm{ncdm}}=(\mnu/3,\mnu/3,\mnu/3)$ for three degenerate species. The identity between the sampled and reported mass is pure bookkeeping. Passing the species masses with seventeen significant digits, not eight decimal places, makes the numerical mapping faithful to the intended masses.

Empirical quantiles are computed from the inverse weighted empirical CDF, with integer holding times treated exactly. Such quantiles do not depend on how identical samples are split into rows. Convergence is assessed with independent chains using rank-normalized split $\hat R$, bulk and tail effective sample sizes, and the Monte Carlo standard error (MCSE) of the 95th-percentile, against thresholds fixed in advance \cite{gelman_rubin1992,vehtari2021_rhat}. A sampled parameter passes when $\hat R\le1.01$; the bulk, tail and quantile ESS are all at least 400; the quantile MCSE is positive and within its limit; and the chronological drift of the quantile and its sensitivity to equal-length chain truncation are within fixed multiples of that limit. For $\mnu$ the MCSE limit applies to the 95th percentile and is absolute: $0.005$\,eV in the first SPT+DESI campaign and $0.001$\,eV in the full CMB+DESI campaign and in all later SPT+DESI campaigns. For other parameters it is $5\%$ of the empirical posterior standard deviation. A group of chains qualifies only if every sampled parameter passes. Chains from different temperatures, solver versions or numerical settings are never pooled. These thresholds are empirical checks, not convergence proofs. History-dependent proposal adaptation is in particular not covered by any theorem we invoke.

In the 28 default-prior chain families examined for this purpose, the largest saved post-burn-in mass is about $0.36$\,eV, far from the prior box $\mnu\le5$\,eV. (Some later trial cohorts with modified numerical settings reached the box during unconverged start-up phases.) Lack of contact with the upper gate does not show that a wider prior would leave the 95th percentile unchanged, as \cref{prop:upper-gate} makes precise.
